\setgebiet{9}{Topologie}{tief}
\chapter{Topologie}
\marginnote{\small\itshape\color{grau}\textbf{Sie brauchen:} Band 8, Kapitel 1--14. Dazu Grenzwerte und Stetigkeit (Gebiet 3), ein wenig Gruppentheorie und lineare Algebra über $\mathbb Z$ (Gebiet 2) für die Kapitel 8, 10 und 11.}[12pt]
\worum{Topologie übt man an Beispielen und Gegenbeispielen: Welche Mengen sind offen, welche Räume homöomorph, welche kompakt, welche zusammenhängend --- und woran erkennt man das? Die zweite Hälfte rechnet: Fundamentalgruppen, Kantenwörter, Homologie, Knotenfärbungen, Barcodes. Die Stolpersteine sind die Stellen, an denen die Anschauung zuverlässig in die Irre führt.}

\section{Was Sie hier üben}
\begin{center}
{\footnotesize\sffamily\setlength{\tabcolsep}{4pt}
\renewcommand{\arraystretch}{1.5}
\begin{tabular}{@{}lP{4.9cm}lcll@{}}
\toprule
\textbf{Nr.} & \textbf{Titel} & \textbf{Sorte} & \textbf{Grad} & \textbf{Zeit} & \textbf{Band/Kap.}\\
\midrule
8.1 & Warum es nur fünf platonische Körper gibt & \typtypisch & \grad2 & 15 min & 8/1\\
8.2 & Die Brücken von Königsberg & \typwerk & \grad1 & 10 min & 8/1\\
8.3 & Der Torus hat $\chi = 0$ & \typwerk & \grad1 & 10 min & 8/1\\
8.4 & Drei ungewohnte Metriken & \typtypisch & \grad2 & 15 min & 8/2\\
8.5 & Die Maximumsmetrik & \typwerk & \grad1 & 10 min & 8/2\\
8.6 & $x = \cos x$ mit Banach & \typwerk & \grad2 & 15 min & 8/2\\
8.7 & Vollständigkeit ist nicht topologisch & \typstolper & \grad2 & 10 min & 8/2\\
8.8 & Die Sorgenfrey-Gerade & \typwerk & \grad2 & 15 min & 8/3\\
8.9 & Inneres, Abschluss, Rand & \typwerk & \grad1 & 10 min & 8/3\\
8.10 & Unendlich viele offene Mengen & \typstolper & \grad1 & 10 min & 8/3\\
8.11 & Homöomorphismen hinschreiben & \typtypisch & \grad2 & 15 min & 8/4\\
8.12 & Das Möbiusband & \typwerk & \grad2 & 15 min & 8/4\\
8.13 & Stetig und bijektiv, aber kein Homöomorphismus & \typstolper & \grad2 & 10 min & 8/4\\
8.14 & Räume durch Zusammenhang unterscheiden & \typtypisch & \grad2 & 15 min & 8/5\\
8.15 & Der Sinuskamm & \typwerk & \grad3 & 20 min & 8/5\\
8.16 & Kompakt oder nicht? & \typtypisch & \grad2 & 15 min & 8/6\\
8.17 & Die Cantor-Menge & \typwerk & \grad2 & 20 min & 8/6\\
8.18 & Die Einheitskugel in $\ell^{2}$ & \typstolper & \grad2 & 15 min & 8/6\\
8.19 & Trennungsaxiome prüfen & \typwerk & \grad2 & 15 min & 8/7\\
8.20 & Eine Folge mit allen Grenzwerten & \typstolper & \grad1 & 10 min & 8/7\\
8.21 & Deformationsretrakte & \typtypisch & \grad2 & 15 min & 8/8\\
8.22 & Der punktierte Torus & \typwerk & \grad2 & 15 min & 8/8\\
8.23 & Überlagerungen und Windungszahl & \typtypisch & \grad3 & 20 min & 8/8\\
8.24 & Brouwer anwenden & \typtypisch & \grad2 & 15 min & 8/9\\
8.25 & Der Schinkensandwich-Satz & \typwerk & \grad3 & 20 min & 8/9\\
8.26 & Fixpunkte ohne Kompaktheit & \typstolper & \grad1 & 10 min & 8/9\\
8.27 & Kantenwörter lesen & \typtypisch & \grad2 & 15 min & 8/10\\
8.28 & Das Wort $aba b^{-1}$ & \typwerk & \grad2 & 15 min & 8/10\\
8.29 & Homologie des Tetraederrands & \typwerk & \grad2 & 20 min & 8/11\\
8.30 & Betti-Zahlen bestimmen & \typtypisch & \grad2 & 20 min & 8/11\\
8.31 & Dreifärbbarkeit & \typwerk & \grad2 & 15 min & 8/12\\
8.32 & Der Achterknoten & \typtypisch & \grad3 & 20 min & 8/12\\
8.33 & Höhenfunktionen auf dem Torus & \typwerk & \grad2 & 20 min & 8/13\\
8.34 & Ein Barcode von Hand & \typwerk & \grad2 & 20 min & 8/14\\
\bottomrule
\end{tabular}}
\end{center}

\section{Woran Sie merken, dass Sie es können}
\begin{itemize}[leftmargin=1.2em,itemsep=2pt]
\item Sie können für eine konkrete Menge in einem konkreten Raum entscheiden, ob sie offen, abgeschlossen, kompakt, zusammenhängend ist --- und den Beweis in zwei Sätzen führen.
\item Sie weisen nach, dass zwei Räume homöomorph sind (Abbildung hinschreiben) oder nicht (eine Invariante finden, die sie trennt).
\item Sie berechnen $\chi$, $\pi_1$ und die Betti-Zahlen einfacher Räume aus einem Modell: Quadrat mit Verklebung, Kantenwort, Triangulierung.
\item Sie kennen die Gegenbeispiele --- Sinuskamm, Cantor-Menge, kofinite Topologie, Einheitskugel in $\ell^{2}$ --- und wissen, welchen falschen Satz jedes von ihnen widerlegt.
\end{itemize}

\section{Aufgaben}

%%%%%%%%%%% Kapitel 1
\stupser{tief}{Zählen Sie Kanten zweimal: jede gehört zu $2$ Flächen und zu $2$ Ecken. Mit $p$-Ecken und $q$ Kanten pro Ecke: $pF = 2K = qE$.}{22pt}
\begin{karte}{tief}{8.1}{\typtypisch}{\meta{2}{15 min}{Band 8, Kap.\ 1}}
\kartentitel{Warum es nur fünf platonische Körper gibt}
\worum{Eulers Polyedersatz $E - K + F = 2$ ist mehr als eine Zählformel: Er schließt Dinge aus. Die klassische Anwendung ist der Beweis, dass es genau fünf regelmäßige konvexe Polyeder gibt.}
\methoden{Euler-Charakteristik, Doppeltes Abzählen, Fallunterscheidung}
Ein regelmäßiges Polyeder habe lauter $p$-Ecke als Flächen, und an jeder Ecke stoßen $q$ Kanten zusammen ($p, q\ge3$).
\begin{enumerate}[label=(\alph*),leftmargin=1.6em,itemsep=2pt]
\item Zeigen Sie $pF = 2K$ und $qE = 2K$.
\item Setzen Sie in $E - K + F = 2$ ein und leiten Sie $\frac1p + \frac1q = \frac12 + \frac1K$ her.
\item Folgern Sie $\frac1p + \frac1q > \frac12$ und bestimmen Sie alle Lösungen $(p, q)$ mit $p, q\ge3$.
\item Berechnen Sie für jede Lösung $K$, $E$, $F$ und benennen Sie den Körper.
\end{enumerate}
\end{karte}

\stupser{gruen}{Zählen Sie an jedem Ufer und jeder Insel die Brücken. Ein Weg, der jede Brücke genau einmal benutzt, betritt und verlässt jeden Zwischenpunkt paarweise.}{22pt}
\begin{karte}{gruen}{8.2}{\typwerk}{\meta{1}{10 min}{Band 8, Kap.\ 1, Aufgabe 1.1}}
\kartentitel{Die Brücken von Königsberg}
\worum{Vier Landteile $A, B, C, D$, sieben Brücken: $A$--$C$ zweimal, $A$--$D$, $B$--$C$ zweimal, $B$--$D$, $C$--$D$. Gesucht ist ein Spaziergang, der jede Brücke genau einmal überquert. Euler zeigte $1736$, dass es keinen gibt --- der erste topologische Beweis.}
\methoden{Gradzählung, Euler-Weg, Invariante}
\begin{enumerate}[label=(\alph*),leftmargin=1.6em,itemsep=2pt]
\item Bestimmen Sie den Grad (die Brückenzahl) jedes Landteils und prüfen Sie: Die Summe ist das Doppelte der Brückenzahl.
\item Beweisen Sie: Hat ein zusammenhängender Graph einen Weg, der jede Kante genau einmal benutzt, so haben höchstens zwei Knoten ungeraden Grad. Folgern Sie, dass es in Königsberg keinen solchen Weg gibt.
\item Welche Brücke müsste man abreißen (oder hinzubauen), damit ein solcher Weg existiert? Und was ist nötig, damit der Weg zum Ausgangspunkt zurückkehrt?
\end{enumerate}
\end{karte}

\stupser{gruen}{Quadrat mit verklebten Gegenseiten: Wie viele Ecken bleiben nach dem Verkleben übrig? Alle vier werden eine.}{22pt}
\begin{karte}{gruen}{8.3}{\typwerk}{\meta{1}{10 min}{Band 8, Kap.\ 1, Aufgabe 1.2}}
\kartentitel{Der Torus hat $\chi = 0$}
\worum{Die Euler-Charakteristik hängt nicht von der Zerlegung ab --- das ist ein Satz, den man an Beispielen prüfen sollte, bevor man ihm glaubt.}
\methoden{Euler-Charakteristik, Quadratmodell, Zählen}
\begin{enumerate}[label=(\alph*),leftmargin=1.6em,itemsep=2pt]
\item Quadratmodell: Der Torus entsteht aus $[0, 1]^{2}$ durch Verkleben von $(0, y)\sim(1, y)$ und $(x, 0)\sim(x, 1)$. Zählen Sie Ecken, Kanten und Flächen \emph{nach} dem Verkleben.
\item Zweite Zerlegung: Teilen Sie das Quadrat in ein $3\times3$-Gitter aus neun kleinen Quadraten und zählen Sie erneut nach dem Verkleben.
\item Dritte Zerlegung (frei): Zerlegen Sie den Torus als \enquote{Ring aus $n$ Vierecken}, also einen Zylinder aus $n$ Rechtecken, dessen Enden verklebt werden. Was ergibt sich für beliebiges $n$?
\end{enumerate}
\end{karte}

%%%%%%%%%%% Kapitel 2
\stupser{tief}{(b) Zwei Punkte, die nicht mit dem Ursprung auf einer Geraden liegen, erreicht man nur über Paris --- den Ursprung.}{22pt}
\begin{karte}{tief}{8.4}{\typtypisch}{\meta{2}{15 min}{Band 8, Kap.\ 2}}
\kartentitel{Drei ungewohnte Metriken}
\worum{Metrik heißt: nichtnegativ und definit, symmetrisch, Dreiecksungleichung. Die Beispiele zeigen, wie verschieden \enquote{Abstand} aussehen kann --- und dass die offenen Kugeln nicht rund sein müssen.}
\methoden{Metrik prüfen, Dreiecksungleichung, Offene Kugeln, Fallunterscheidung}
\begin{enumerate}[label=(\alph*),leftmargin=1.6em,itemsep=2pt]
\item Die diskrete Metrik auf einer beliebigen Menge: $d(x, y) = 0$ für $x = y$, sonst $1$. Prüfen Sie die Axiome. Was sind die offenen Kugeln $B_r(x)$ für $r\le1$ und $r > 1$? Welche Mengen sind offen?
\item Die französische Eisenbahnmetrik auf $\mathbb R^{2}$: $d(x, y) = |x - y|$, falls $x, y$ und $0$ auf einer Geraden liegen, sonst $d(x, y) = |x| + |y|$. Prüfen Sie die Dreiecksungleichung (Fallunterscheidung) und skizzieren Sie $B_r(x)$ für $r < |x|$ und für $r > |x|$.
\item Auf $\mathbb R$ sei $d(x, y) = |\arctan x - \arctan y|$. Warum ist das eine Metrik? Welche Teilmengen von $\mathbb R$ sind offen? Ist $(\mathbb R, d)$ vollständig? (Betrachten Sie $x_n = n$.)
\end{enumerate}
\end{karte}

\stupser{gruen}{Dreiecksungleichung: Das Maximum über $i$ von $|x_i - z_i|\le|x_i - y_i| + |y_i - z_i|$ ist höchstens die Summe der beiden Maxima.}{22pt}
\begin{karte}{gruen}{8.5}{\typwerk}{\meta{1}{10 min}{Band 8, Kap.\ 2, Aufgabe 2.1}}
\kartentitel{Die Maximumsmetrik}
\worum{Auf $\mathbb R^{n}$ gibt es viele Metriken. $d_\infty(x, y) = \max_i|x_i - y_i|$ ist die bequemste für Abschätzungen --- und erzeugt dieselbe Topologie wie die euklidische.}
\methoden{Metrik prüfen, Dreiecksungleichung, Offene Kugeln, Äquivalente Metriken}
\begin{enumerate}[label=(\alph*),leftmargin=1.6em,itemsep=2pt]
\item Zeigen Sie, dass $d_\infty$ eine Metrik auf $\mathbb R^{n}$ ist.
\item Zeichnen Sie die Einheitskugeln $\{x\colon d(x, 0)\le1\}$ in $\mathbb R^{2}$ für $d_1(x, y) = \sum|x_i - y_i|$, $d_2$ (euklidisch) und $d_\infty$.
\item Zeigen Sie $d_\infty\le d_2\le d_1\le n\,d_\infty$ und folgern Sie, dass alle drei dieselben offenen Mengen haben.
\end{enumerate}
\end{karte}

\stupser{gruen}{Auf $[0{,}5;\ 1]$ ist $|\cos x - \cos y| = |\sin\xi|\,|x - y|\le\sin(1)\,|x - y|$ mit $\sin1 = 0{,}8415$.}{22pt}
\begin{karte}{gruen}{8.6}{\typwerk}{\meta{2}{15 min}{Band 8, Kap.\ 2, Aufgabe 2.2}}
\kartentitel{$x = \cos x$ mit Banach}
\worum{Der Banachsche Fixpunktsatz liefert Existenz, Eindeutigkeit, ein Verfahren und eine Fehlerschranke auf einmal. Die Gleichung $x = \cos x$ ist das Musterbeispiel.}
\methoden{Banachscher Fixpunktsatz, Fixpunktiteration, Mittelwertsatz, Fehlerschranke}
\begin{enumerate}[label=(\alph*),leftmargin=1.6em,itemsep=2pt]
\item Finden Sie ein abgeschlossenes Intervall $X$ mit $\cos(X)\subseteq X$.
\item Zeigen Sie mit dem Mittelwertsatz, dass $\cos$ auf $X$ eine Kontraktion ist, und bestimmen Sie die Konstante $q$.
\item Iterieren Sie von $x_0 = 0{,}75$ aus drei Schritte und geben Sie den Fixpunkt auf sechs Stellen an.
\item A-priori-Schranke: $d(x_n, x^{*})\le\frac{q^{n}}{1 - q}\,d(x_1, x_0)$. Wie viele Schritte garantieren einen Fehler unter $10^{-6}$? Vergleichen Sie mit dem tatsächlichen Fehler nach $10$ Schritten.
\end{enumerate}
\end{karte}

\stupser{akzent}{$x\mapsto\tan\bigl(\pi(x - \frac12)\bigr)$ ist ein Homöomorphismus $(0, 1)\to\mathbb R$. Links ist $\frac1n$ eine Cauchy-Folge ohne Grenzwert, rechts gibt es so etwas nicht.}{22pt}
\begin{karte}{akzent}{8.7}{\typstolper}{\meta{2}{10 min}{Band 8, Kap.\ 2}}
\kartentitel{Vollständigkeit ist nicht topologisch}
\worum{Vollständigkeit, Beschränktheit, Cauchy-Folge --- diese Begriffe gehören zur Metrik, nicht zur Topologie. Homöomorphe Räume können sich in ihnen unterscheiden. Wer das vergisst, \enquote{beweist} falsche Sätze.}
\methoden{Gegenbeispiel, Homöomorphismus, Cauchy-Folge, Invariante}
\begin{enumerate}[label=(\alph*),leftmargin=1.6em,itemsep=2pt]
\item Zeigen Sie, dass $(0, 1)$ und $\mathbb R$ homöomorph sind, $(0, 1)$ aber nicht vollständig und $\mathbb R$ vollständig ist.
\item Was ist die Folge $\frac1n$ in $(0, 1)$ unter dem Homöomorphismus? Ist das Bild eine Cauchy-Folge?
\item Beschränktheit: Ist \enquote{beschränkt} eine topologische Eigenschaft? Und \enquote{kompakt}? Begründen Sie den Unterschied.
\item Welche der folgenden Aussagen sind richtig? (i) Ein stetiges Bild einer Cauchy-Folge ist eine Cauchy-Folge. (ii) Ein gleichmäßig stetiges Bild einer Cauchy-Folge ist eine Cauchy-Folge. (iii) Ein Homöomorphismus bildet konvergente Folgen auf konvergente Folgen ab.
\end{enumerate}
\end{karte}

%%%%%%%%%%% Kapitel 3
\stupser{gruen}{(d) Zu jedem $x$ gibt es ein Basiselement $B$ mit $x\in B\subseteq[x, x + 1)$ --- und dann ist $x = \min B$. Verschiedene $x$, verschiedene $B$.}{22pt}
\begin{karte}{gruen}{8.8}{\typwerk}{\meta{2}{15 min}{Band 8, Kap.\ 3, Aufgabe 3.1}}
\kartentitel{Die Sorgenfrey-Gerade}
\worum{$\mathbb R_\ell$ ist $\mathbb R$ mit der von den halboffenen Intervallen $[a, b)$ erzeugten Topologie. Sie ist das Standardgegenbeispiel der mengentheoretischen Topologie: harmlos aussehend, mit überraschenden Eigenschaften.}
\methoden{Basiskriterium, Offen und abgeschlossen, Feinere Topologie, Abzählbarkeitsargument}
\begin{enumerate}[label=(\alph*),leftmargin=1.6em,itemsep=2pt]
\item Zeigen Sie, dass $\mathcal B = \{[a, b)\colon a < b\}$ eine Basis ist (Durchschnitt zweier Basiselemente ist wieder eines oder leer).
\item Ist $[0, 1)$ offen? Abgeschlossen? Und $(0, 1)$? Und $[0, 1]$?
\item Zeigen Sie: Jede in $\mathbb R$ offene Menge ist in $\mathbb R_\ell$ offen, aber nicht umgekehrt. Welche Topologie ist feiner?
\item Zeigen Sie, dass $\mathbb R_\ell$ keine abzählbare Basis hat. (Warum hat trotzdem jeder Punkt eine abzählbare Umgebungsbasis?)
\end{enumerate}
\end{karte}

\stupser{gruen}{Dreimal dieselbe Menge, drei Räume. In $\mathbb R$: Jedes Intervall enthält Irrationales (kein Inneres) und Rationales (Abschluss $[0, 1]$).}{22pt}
\begin{karte}{gruen}{8.9}{\typwerk}{\meta{1}{10 min}{Band 8, Kap.\ 3, Aufgabe 3.2}}
\kartentitel{Inneres, Abschluss, Rand}
\worum{Inneres, Abschluss und Rand hängen vom umgebenden Raum ab, nicht nur von der Menge. Die Aufgabe macht das an $A = \mathbb Q\cap[0, 1]$ sichtbar.}
\methoden{Inneres und Abschluss, Rand, Teilraumtopologie, Dichtheit}
Bestimmen Sie $\operatorname{int}A$, $\overline A$ und $\partial A$ für $A = \mathbb Q\cap[0, 1]$
\begin{enumerate}[label=(\alph*),leftmargin=1.6em,itemsep=2pt]
\item in $\mathbb R$ mit der Standardtopologie,
\item in $\mathbb Q$ mit der Teilraumtopologie,
\item in $\mathbb R$ mit der diskreten Topologie.
\end{enumerate}
Prüfen Sie jeweils $\partial A = \overline A\setminus\operatorname{int}A$.
\end{karte}

\stupser{akzent}{$\bigcap_n(-\frac1n, \frac1n) = \{0\}$: Unendlich viele offene Mengen, ein Punkt. Die Axiome erlauben nur \emph{endliche} Durchschnitte --- und das ist kein Schönheitsfehler.}{22pt}
\begin{karte}{akzent}{8.10}{\typstolper}{\meta{1}{10 min}{Band 8, Kap.\ 3}}
\kartentitel{Unendlich viele offene Mengen}
\worum{Beliebige Vereinigungen offener Mengen sind offen, aber nur endliche Durchschnitte. Analog vertauscht der Abschluss mit endlichen, nicht mit unendlichen Vereinigungen. Die Gegenbeispiele sind klein und sollten im Kopf sein.}
\methoden{Gegenbeispiel, Offene Mengen, Inneres und Abschluss}
\begin{enumerate}[label=(\alph*),leftmargin=1.6em,itemsep=2pt]
\item Geben Sie unendlich viele offene Mengen in $\mathbb R$ an, deren Durchschnitt nicht offen ist, und unendlich viele abgeschlossene, deren Vereinigung nicht abgeschlossen ist.
\item Zeigen Sie $\overline{A\cup B} = \overline A\cup\overline B$ und geben Sie Mengen $A_n$ mit $\overline{\bigcup A_n}\neq\bigcup\overline{A_n}$ an.
\item Zeigen Sie $\partial(A\cup B)\subseteq\partial A\cup\partial B$ und finden Sie $A, B$ mit echter Inklusion.
\item Was wäre, wenn man beliebige Durchschnitte offener Mengen zuließe? Zeigen Sie: In $\mathbb R$ wäre dann jede Menge offen.
\end{enumerate}
\end{karte}

%%%%%%%%%%% Kapitel 4
\stupser{tief}{Ein Homöomorphismus ist eine stetige Bijektion mit stetiger Umkehrung. Am schnellsten: beide Abbildungen hinschreiben und ihre Stetigkeit ablesen.}{22pt}
\begin{karte}{tief}{8.11}{\typtypisch}{\meta{2}{15 min}{Band 8, Kap.\ 4}}
\kartentitel{Homöomorphismen hinschreiben}
\worum{\enquote{Topologisch gleich} beweist man konstruktiv: Man gibt die Abbildung an. Die Aufgabe sammelt die Standardfälle, die man auswendig können sollte.}
\methoden{Homöomorphismus, Stetigkeit, Umkehrabbildung}
Geben Sie jeweils einen Homöomorphismus mit Umkehrabbildung an:
\begin{enumerate}[label=(\alph*),leftmargin=1.6em,itemsep=2pt]
\item $(0, 1)\to\mathbb R$ und $(a, b)\to(c, d)$.
\item Die offene Einheitskreisscheibe $D = \{x\in\mathbb R^{2}\colon|x| < 1\}\to\mathbb R^{2}$.
\item Der Rand des Quadrats $[-1, 1]^{2}$ $\to$ die Kreislinie $S^{1}$.
\item $S^{1}\setminus\{(0, 1)\}\to\mathbb R$ (stereographische Projektion). Folgern Sie: $S^{1}$ ist die Einpunktkompaktifizierung von $\mathbb R$.
\item $\mathbb R^{n}\setminus\{0\}\to S^{n-1}\times(0, \infty)$.
\end{enumerate}
\end{karte}

\stupser{gruen}{Modell: $[0, 1]^{2}$ mit $(0, y)\sim(1, 1 - y)$. Verfolgen Sie die Randkante $y = 0$: Sie läuft bei $x = 1$ weiter als $y = 1$.}{22pt}
\begin{karte}{gruen}{8.12}{\typwerk}{\meta{2}{15 min}{Band 8, Kap.\ 4, Aufgabe 4.1}}
\kartentitel{Das Möbiusband}
\worum{Das Möbiusband ist der erste Raum, an dem die Anschauung versagt: ein Rand, eine Seite, und Schnitte, die nicht trennen. Papier und Schere sind erlaubt --- aber die Antworten sollen am Modell begründet werden.}
\methoden{Quotientenraum, Verklebung verfolgen, Euler-Charakteristik, Orientierbarkeit}
\begin{enumerate}[label=(\alph*),leftmargin=1.6em,itemsep=2pt]
\item Wie viele Randkomponenten hat das Möbiusband? Begründen Sie am Quadratmodell.
\item Was entsteht beim Schnitt längs der Mittellinie $y = \frac12$? Wie viele Verdrehungen hat das Ergebnis, und ist es einseitig?
\item Was entsteht beim Schnitt längs $y = \frac13$?
\item Bestimmen Sie $\chi$ über eine Zerlegung des Quadrats. Vergleichen Sie mit Zylinder und Torus.
\end{enumerate}
\end{karte}

\stupser{akzent}{$f\colon[0, 1)\to S^{1}$, $t\mapsto e^{2\pi it}$, ist stetig und bijektiv. Die Umkehrung ist bei $(1, 0)$ unstetig --- oder: $[0, 1)$ ist nicht kompakt, $S^{1}$ schon.}{22pt}
\begin{karte}{akzent}{8.13}{\typstolper}{\meta{2}{10 min}{Band 8, Kap.\ 4}}
\kartentitel{Stetig und bijektiv, aber kein Homöomorphismus}
\worum{In der linearen Algebra ist eine bijektive lineare Abbildung automatisch ein Isomorphismus. In der Topologie ist eine stetige Bijektion \emph{nicht} automatisch ein Homöomorphismus --- außer unter einer Zusatzvoraussetzung, die man kennen sollte.}
\methoden{Gegenbeispiel, Homöomorphismus, Kompaktheit, Umkehrabbildung}
\begin{enumerate}[label=(\alph*),leftmargin=1.6em,itemsep=2pt]
\item Zeigen Sie, dass $f\colon[0, 1)\to S^{1}$, $t\mapsto(\cos2\pi t, \sin2\pi t)$, stetig und bijektiv ist, aber $f^{-1}$ nicht stetig.
\item Zeigen Sie dasselbe für $\operatorname{id}\colon(\mathbb R, \text{diskret})\to(\mathbb R, \text{Standard})$.
\item Beweisen Sie: Eine stetige Bijektion von einem kompakten Raum auf einen Hausdorff-Raum ist ein Homöomorphismus. Welche Voraussetzung ist in (a) und (b) verletzt?
\item Warum folgt aus (c), dass es keine stetige Bijektion $[0, 1]\to S^{1}$ gibt, obwohl beide kompakt sind? (Hinweis: Entfernen Sie einen Punkt.)
\end{enumerate}
\end{karte}

%%%%%%%%%%% Kapitel 5
\stupser{tief}{Zusammenhang ist eine Invariante --- und bleibt es, wenn man aus beiden Räumen \enquote{entsprechende} Punkte entfernt. Zählen Sie, wie viele Teile ein Raum nach Entfernen eines Punkts hat.}{22pt}
\begin{karte}{tief}{8.14}{\typtypisch}{\meta{2}{15 min}{Band 8, Kap.\ 5}}
\kartentitel{Räume durch Zusammenhang unterscheiden}
\worum{Wie beweist man, dass zwei Räume \emph{nicht} homöomorph sind? Man braucht eine Invariante, in der sie sich unterscheiden. Zusammenhang --- vor und nach dem Entfernen von Punkten --- ist die einfachste.}
\methoden{Zusammenhang, Punkte entfernen, Invariante, Zwischenwertsatz}
\begin{enumerate}[label=(\alph*),leftmargin=1.6em,itemsep=2pt]
\item Zeigen Sie: $[0, 1]$ und $S^{1}$ sind nicht homöomorph. Ebenso $[0, 1]$ und $[0, 1)$.
\item Zeigen Sie: $\mathbb R$ und $\mathbb R^{2}$ sind nicht homöomorph.
\item Die Buchstaben I, L, T, X, O, P als Teilmengen der Ebene: Welche sind untereinander homöomorph? Begründen Sie mit der Zahl der Punkte, deren Entfernung den Buchstaben in $1$, $2$, $3$ oder $4$ Teile zerlegt.
\item Beweisen Sie mit dem Zwischenwertsatz: Zu jeder stetigen Funktion $f\colon S^{1}\to\mathbb R$ gibt es $x$ mit $f(x) = f(-x)$. (Borsuk--Ulam in Dimension $1$.)
\end{enumerate}
\end{karte}

\stupser{gruen}{Zusammenhängend: $S$ ist der Abschluss des Graphen, und der Abschluss einer zusammenhängenden Menge ist zusammenhängend. Wegzusammenhängend: Ein Weg aus $A$ heraus müsste die Oszillationen $\pm1$ überleben.}{22pt}
\begin{karte}{gruen}{8.15}{\typwerk}{\meta{3}{20 min}{Band 8, Kap.\ 5, Aufgabe 5.1}}
\kartentitel{Der Sinuskamm}
\worum{$S = G\cup A$ mit $G = \{(x, \sin\frac1x)\colon0 < x\le1\}$ und $A = \{0\}\times[-1, 1]$. Das Standardbeispiel dafür, dass Zusammenhang und Wegzusammenhang verschieden sind.}
\methoden{Zusammenhang, Wegzusammenhang, Abschluss, Widerspruchsbeweis, Grenzwerte}
\begin{enumerate}[label=(\alph*),leftmargin=1.6em,itemsep=2pt]
\item Zeigen Sie: $G$ ist wegzusammenhängend, und $\overline G = S$.
\item Folgern Sie, dass $S$ zusammenhängend ist.
\item Zeigen Sie, dass $S$ nicht wegzusammenhängend ist: Angenommen $\gamma\colon[0, 1]\to S$ ist ein Weg mit $\gamma(0)\in A$ und $\gamma(1)\in G$. Sei $t_0 = \sup\{t\colon\gamma(t)\in A\}$. Leiten Sie aus der Stetigkeit bei $t_0$ einen Widerspruch her.
\item Was ändert sich, wenn man $A$ durch den einzelnen Punkt $(0, 0)$ ersetzt? Durch die Strecke $\{0\}\times[-2, 2]$?
\end{enumerate}
\end{karte}

%%%%%%%%%%% Kapitel 6
\stupser{tief}{In $\mathbb R^{n}$: kompakt $\iff$ abgeschlossen und beschränkt (Heine--Borel). Für \enquote{nicht kompakt} genügt eine offene Überdeckung ohne endliche Teilüberdeckung --- oder eine Folge ohne Häufungspunkt.}{22pt}
\begin{karte}{tief}{8.16}{\typtypisch}{\meta{2}{15 min}{Band 8, Kap.\ 6}}
\kartentitel{Kompakt oder nicht?}
\worum{Kompaktheit ist der Begriff, der den Satz vom Maximum, die gleichmäßige Stetigkeit und Bolzano--Weierstraß trägt. Man muss sie an konkreten Mengen schnell entscheiden können.}
\methoden{Kompaktheit, Heine-Borel, Offene Überdeckung, Gegenbeispiel}
\begin{enumerate}[label=(\alph*),leftmargin=1.6em,itemsep=2pt]
\item Geben Sie für $(0, 1)$ eine offene Überdeckung ohne endliche Teilüberdeckung an; ebenso für $\mathbb Z\subseteq\mathbb R$ und für $[0, \infty)$.
\item Welche der folgenden Mengen sind kompakt: $\{0\}\cup\{\frac1n\colon n\in\mathbb N\}$; $\{\frac1n\}$ ohne die $0$; die Einheitssphäre $S^{n-1}$; $\{(x, y)\colon xy = 1\}$; $\{(x, y)\colon x^{2} + y^{2}\le1,\ y > 0\}$; $\mathbb Q\cap[0, 1]$ (in $\mathbb R$)?
\item Zeigen Sie: Das stetige Bild einer kompakten Menge ist kompakt. Folgern Sie den Satz vom Maximum.
\item Zeigen Sie: Eine kompakte Teilmenge eines metrischen Raums ist abgeschlossen und beschränkt. (Die Umkehrung gilt nicht --- Aufgabe 8.18.)
\end{enumerate}
\end{karte}

\stupser{gruen}{$C = \bigcap C_n$, jedes $C_n$ eine endliche Vereinigung abgeschlossener Intervalle. Elemente von $C$: Ternärbrüche nur mit Ziffern $0$ und $2$.}{22pt}
\begin{karte}{gruen}{8.17}{\typwerk}{\meta{2}{20 min}{Band 8, Kap.\ 6, Aufgabe 6.1}}
\kartentitel{Die Cantor-Menge}
\worum{Aus $[0, 1]$ entfernt man das offene mittlere Drittel, aus den beiden Resten wieder, und so fort. Was übrig bleibt, ist das vielseitigste Gegenbeispiel der Analysis: kompakt, überabzählbar, vom Maß null, ohne isolierte Punkte, total unzusammenhängend.}
\methoden{Kompaktheit, Ternärdarstellung, Cantors Diagonalargument, Geometrische Reihe, Zusammenhang}
\begin{enumerate}[label=(\alph*),leftmargin=1.6em,itemsep=2pt]
\item Zeigen Sie, dass $C$ kompakt ist.
\item Zeigen Sie: $x\in C$ genau dann, wenn $x$ eine Ternärdarstellung ohne die Ziffer $1$ hat. Folgern Sie, dass $C$ überabzählbar ist.
\item Berechnen Sie die Gesamtlänge der entfernten Intervalle. Was folgt für das \enquote{Maß} von $C$?
\item Zeigen Sie: $C$ enthält kein Intervall (total unzusammenhängend), aber keinen isolierten Punkt (perfekt).
\item Ist $\frac14$ in $C$? Und $\frac34$? Und $\frac1{10}$?
\end{enumerate}
\end{karte}

\stupser{akzent}{Die Einheitsvektoren $e_n$ haben paarweise Abstand $\sqrt2$. Eine Folge ohne konvergente Teilfolge in einer abgeschlossenen, beschränkten Menge.}{22pt}
\begin{karte}{akzent}{8.18}{\typstolper}{\meta{2}{15 min}{Band 8, Kap.\ 6}}
\kartentitel{Die Einheitskugel in $\ell^{2}$}
\worum{\enquote{Abgeschlossen und beschränkt, also kompakt} --- in $\mathbb R^{n}$ richtig, in unendlichdimensionalen Räumen falsch. Der Fehler ist häufig, weil Heine--Borel so vertraut ist.}
\methoden{Gegenbeispiel, Kompaktheit, Folgenkompaktheit, Offene Überdeckung}
Sei $\ell^{2}$ der Raum der quadratsummierbaren Folgen mit $\|x\| = \bigl(\sum x_k^{2}\bigr)^{1/2}$, und $\overline B = \{x\colon\|x\|\le1\}$.
\begin{enumerate}[label=(\alph*),leftmargin=1.6em,itemsep=2pt]
\item Zeigen Sie, dass $\overline B$ abgeschlossen und beschränkt ist.
\item Sei $e_n$ die Folge mit $1$ an der Stelle $n$ und $0$ sonst. Berechnen Sie $\|e_n - e_m\|$ für $n\neq m$ und folgern Sie, dass $(e_n)$ keine konvergente Teilfolge hat.
\item Geben Sie eine offene Überdeckung von $\overline B$ ohne endliche Teilüberdeckung an. (Kugeln vom Radius $\frac12$ um alle Punkte --- wie viele $e_n$ liegen in einer?)
\item Wo genau geht der Beweis von Heine--Borel (Intervallschachtelung, Halbieren der Würfel) in unendlicher Dimension schief?
\end{enumerate}
\end{karte}

%%%%%%%%%%% Kapitel 7
\stupser{gruen}{(a) Zwei nichtleere offene Mengen der kofiniten Topologie haben endliche Komplemente --- ihr Durchschnitt ist nie leer.}{22pt}
\begin{karte}{gruen}{8.19}{\typwerk}{\meta{2}{15 min}{Band 8, Kap.\ 7, Aufgabe 7.1}}
\kartentitel{Trennungsaxiome prüfen}
\worum{$T_0$: je zwei Punkte durch eine offene Menge unterscheidbar; $T_1$: Einpunktmengen abgeschlossen; $T_2$ (Hausdorff): disjunkte Umgebungen; $T_4$ (normal): disjunkte abgeschlossene Mengen trennbar. Die Beispiele zeigen, dass jede Stufe eine echte ist.}
\methoden{Trennungsaxiome, Hausdorff-Eigenschaft, Gegenbeispiel, Quotientenraum}
Untersuchen Sie $T_0$, $T_1$, $T_2$ für
\begin{enumerate}[label=(\alph*),leftmargin=1.6em,itemsep=2pt]
\item $\mathbb R$ mit der kofiniten Topologie (offen: $\emptyset$ und alle Mengen mit endlichem Komplement). Ist dieser Raum kompakt?
\item den Sierpiński-Raum $S = \{0, 1\}$ mit $\mathcal T = \{\emptyset, \{1\}, S\}$;
\item den Quotienten $\mathbb R/\mathbb Q$ (mit $x\sim y\iff x - y\in\mathbb Q$);
\item jeden metrischen Raum --- zeigen Sie hier zusätzlich $T_4$ mit der Funktion $f(x) = \frac{d(x, A)}{d(x, A) + d(x, B)}$.
\end{enumerate}
\end{karte}

\stupser{akzent}{In der kofiniten Topologie enthält jede Umgebung jedes Punkts fast alle Glieder jeder injektiven Folge. Die Folge $x_n = n$ konvergiert gegen $\pi$, gegen $0$, gegen alles.}{22pt}
\begin{karte}{akzent}{8.20}{\typstolper}{\meta{1}{10 min}{Band 8, Kap.\ 7}}
\kartentitel{Eine Folge mit allen Grenzwerten}
\worum{\enquote{Der Grenzwert} einer Folge --- der bestimmte Artikel setzt Eindeutigkeit voraus, und die gilt nur in Hausdorff-Räumen. Die Aufgabe zeigt, was ohne $T_2$ passiert.}
\methoden{Gegenbeispiel, Hausdorff-Eigenschaft, Konvergenz, Grenzwerte}
\begin{enumerate}[label=(\alph*),leftmargin=1.6em,itemsep=2pt]
\item Zeigen Sie: In $\mathbb R$ mit der kofiniten Topologie konvergiert jede Folge mit paarweise verschiedenen Gliedern gegen \emph{jeden} Punkt.
\item Gegen welche Punkte konvergiert die konstante Folge $x_n = 5$? Die Folge $1, 2, 1, 2, \dots$?
\item Beweisen Sie: In einem Hausdorff-Raum hat jede Folge höchstens einen Grenzwert.
\item Im Sierpiński-Raum (Aufgabe 8.19b): Wogegen konvergiert die konstante Folge $1, 1, 1, \dots$? Und $0, 0, 0, \dots$?
\end{enumerate}
\end{karte}

%%%%%%%%%%% Kapitel 8
\stupser{tief}{Ein Deformationsretrakt braucht eine Homotopie $H(x, t)$ mit $H(x, 0) = x$, $H(x, 1)\in A$ und $H(a, t) = a$ für $a\in A$. Meist tut es die gerade Verbindungsstrecke: $H(x, t) = (1 - t)x + t\,r(x)$.}{22pt}
\begin{karte}{tief}{8.21}{\typtypisch}{\meta{2}{15 min}{Band 8, Kap.\ 8}}
\kartentitel{Deformationsretrakte}
\worum{Homotopieäquivalente Räume haben dieselbe Fundamentalgruppe. Die wichtigste Quelle solcher Äquivalenzen: Ein Raum zieht sich auf einen Teilraum zusammen. Die Aufgabe übt, die Homotopie hinzuschreiben.}
\methoden{Deformationsretrakt, Homotopie, Fundamentalgruppe}
Zeigen Sie, dass $A$ ein Deformationsretrakt von $X$ ist, und geben Sie $\pi_1(X)$ an:
\begin{enumerate}[label=(\alph*),leftmargin=1.6em,itemsep=2pt]
\item $X = \mathbb R^{2}\setminus\{0\}$, $A = S^{1}$.
\item $X = \{x\in\mathbb R^{2}\colon1\le|x|\le2\}$ (Kreisring), $A = \{|x| = 1\}$.
\item $X = \mathbb R^{3}\setminus z\text{-Achse}$, $A = $ Einheitskreis in der $xy$-Ebene.
\item $X$ = Möbiusband, $A$ = Mittellinie. (Was folgt für $\pi_1$ des Möbiusbands?)
\item $X = \mathbb R^{2}\setminus\{p, q\}$, $A$ = eine \enquote{Acht} aus zwei Kreisen um $p$ und $q$, die sich in einem Punkt berühren. (Skizze genügt; was ist $\pi_1$?)
\end{enumerate}
\end{karte}

\stupser{gruen}{Das Quadrat ohne Mittelpunkt zieht sich radial auf seinen Rand zusammen. Der Rand ist nach dem Verkleben eine Acht: $S^{1}\vee S^{1}$.}{22pt}
\begin{karte}{gruen}{8.22}{\typwerk}{\meta{2}{15 min}{Band 8, Kap.\ 8, Aufgabe 8.1}}
\kartentitel{Der punktierte Torus}
\worum{$\pi_1(T^{2}) = \mathbb Z\times\mathbb Z$ ist abelsch. Entfernt man einen Punkt, wird die Gruppe frei --- die Relation $ab = ba$ lebte in der Fläche, die man weggenommen hat.}
\methoden{Deformationsretrakt, Fundamentalgruppe, Quadratmodell, Freie Gruppe}
\begin{enumerate}[label=(\alph*),leftmargin=1.6em,itemsep=2pt]
\item Zeigen Sie, dass $[0, 1]^{2}\setminus\{(\frac12, \frac12)\}$ sich auf den Rand des Quadrats zusammenzieht.
\item Was ist der Rand des Quadrats nach den Torus-Verklebungen? Folgern Sie $\pi_1(T^{2}\setminus\{p\})$.
\item Vergleichen Sie mit $\pi_1(T^{2}) = \langle a, b\mid aba^{-1}b^{-1}\rangle$: Woher kommt die Relation, und warum verschwindet sie?
\item Zeigen Sie, dass $\pi_1(T^{2}\setminus\{p\})$ nicht abelsch ist (es genügt ein Homomorphismus auf eine nichtabelsche Gruppe, etwa $S_3$).
\end{enumerate}
\end{karte}

\stupser{tief}{(a) Jede Schleife in $S^{1}$ hebt sich eindeutig zu einem Weg in $\mathbb R$ mit Anfang $0$; sein Endpunkt ist eine ganze Zahl --- die Windungszahl.}{22pt}
\begin{karte}{tief}{8.23}{\typtypisch}{\meta{3}{20 min}{Band 8, Kap.\ 8}}
\kartentitel{Überlagerungen und Windungszahl}
\worum{$p\colon\mathbb R\to S^{1}$, $t\mapsto e^{2\pi it}$, ist die Mutter aller Überlagerungen. Mit ihr berechnet man $\pi_1(S^{1})$, den Grad von $z\mapsto z^{n}$ und --- über $S^{2}\to\mathbb{RP}^{2}$ --- die Fundamentalgruppe der projektiven Ebene.}
\methoden{Überlagerung, Hochheben von Wegen, Windungszahl, Fundamentalgruppe}
\begin{enumerate}[label=(\alph*),leftmargin=1.6em,itemsep=2pt]
\item Erklären Sie, wie das Hochheben von Wegen jeder Schleife $\gamma$ in $S^{1}$ (Basispunkt $1$) eine ganze Zahl $\deg\gamma$ zuordnet, und warum homotope Schleifen denselben Grad haben.
\item Bestimmen Sie den Grad von $\gamma_n(t) = e^{2\pi int}$ und von $t\mapsto e^{2\pi i\sin(\pi t)}$.
\item Folgern Sie $\pi_1(S^{1})\cong\mathbb Z$ (Sie dürfen benutzen, dass Schleifen gleichen Grades homotop sind).
\item $\mathbb{RP}^{2} = S^{2}/(x\sim-x)$. Zeigen Sie, dass $S^{2}\to\mathbb{RP}^{2}$ eine zweiblättrige Überlagerung ist, und folgern Sie $\pi_1(\mathbb{RP}^{2})\cong\mathbb Z/2$. Welche Schleife ist der Erzeuger, und warum ist ihr Quadrat trivial?
\end{enumerate}
\end{karte}

%%%%%%%%%%% Kapitel 9
\stupser{tief}{(b) Gäbe es eine Retraktion $r\colon D^{2}\to S^{1}$, so hätte $x\mapsto-r(x)$ keinen Fixpunkt. Gäbe es $f$ ohne Fixpunkt, so definiert der Strahl von $f(x)$ durch $x$ eine Retraktion.}{22pt}
\begin{karte}{tief}{8.24}{\typtypisch}{\meta{2}{15 min}{Band 8, Kap.\ 9}}
\kartentitel{Brouwer anwenden}
\worum{Jede stetige Abbildung $D^{n}\to D^{n}$ hat einen Fixpunkt. Der Satz ist nicht konstruktiv, aber ungemein anwendbar --- und er ist äquivalent zu einer Aussage über Retraktionen, die man kennen sollte.}
\methoden{Brouwerscher Fixpunktsatz, Retraktion, Zwischenwertsatz, Widerspruchsbeweis}
\begin{enumerate}[label=(\alph*),leftmargin=1.6em,itemsep=2pt]
\item Eine Landkarte des Landes liegt irgendwo im Land auf dem Boden. Zeigen Sie: Genau ein Punkt der Karte liegt über dem Punkt, den er darstellt. (Warum \enquote{genau ein}? Welcher Fixpunktsatz reicht hier?)
\item Zeigen Sie: Brouwer für $D^{2}$ ist äquivalent zu \enquote{es gibt keine Retraktion $D^{2}\to S^{1}$}.
\item Beweisen Sie Brouwer in Dimension $1$ ($f\colon[0, 1]\to[0, 1]$) mit dem Zwischenwertsatz.
\item Sei $M$ eine $n\times n$-Matrix mit lauter positiven Einträgen. Zeigen Sie mit Brouwer auf dem Simplex $\Delta = \{x\ge0, \sum x_i = 1\}$, dass $M$ einen Eigenvektor mit positiven Einträgen zu einem positiven Eigenwert hat.
\end{enumerate}
\end{karte}

\stupser{gruen}{Zu jeder Richtung $u\in S^{n-1}$ gibt es eine Hyperebene senkrecht zu $u$, die $A_n$ halbiert. Die Volumina von $A_1, \dots, A_{n-1}$ auf der $u$-Seite geben $f(u)\in\mathbb R^{n-1}$ --- und Borsuk--Ulam liefert $f(u) = f(-u)$.}{22pt}
\begin{karte}{gruen}{8.25}{\typwerk}{\meta{3}{20 min}{Band 8, Kap.\ 9, Aufgabe 9.1}}
\kartentitel{Der Schinkensandwich-Satz}
\worum{Zu $n$ beschränkten messbaren Mengen im $\mathbb R^{n}$ gibt es eine Hyperebene, die alle gleichzeitig halbiert. Der Beweis ist eine Anwendung von Borsuk--Ulam --- und ein Muster dafür, wie man einen topologischen Satz auf ein Maßproblem ansetzt.}
\methoden{Satz von Borsuk-Ulam, Zwischenwertsatz, Stetigkeitsargument, Konstruktion einer Hilfsfunktion}
\begin{enumerate}[label=(\alph*),leftmargin=1.6em,itemsep=2pt]
\item $n = 1$: Zeigen Sie mit dem Zwischenwertsatz, dass es zu einer beschränkten messbaren Menge $A\subseteq\mathbb R$ einen Punkt gibt, der sie halbiert.
\item $n = 2$: Zu jeder Richtung $u\in S^{1}$ gibt es eine Gerade senkrecht zu $u$, die $A_2$ halbiert. Sei $f(u)$ der Flächeninhalt von $A_1$ auf der Seite, in die $u$ zeigt. Zeigen Sie $f(u) + f(-u) = |A_1|$ und wenden Sie Borsuk--Ulam ($S^{1}\to\mathbb R$) an.
\item Formulieren Sie den Beweis für allgemeines $n$ mit $f\colon S^{n-1}\to\mathbb R^{n-1}$.
\item Wo braucht man, dass die halbierende Hyperebene stetig von $u$ abhängt, und wie behebt man das Problem, wenn sie nicht eindeutig ist? (Hinweis: Mittelwert aller halbierenden Positionen.)
\end{enumerate}
\end{karte}

\stupser{akzent}{Offene Scheibe: $f(x) = \frac{x + e_1}2$ schiebt alles Richtung Rand; der Fixpunkt $e_1$ liegt nicht mehr im Raum. Kreisring: eine Drehung.}{22pt}
\begin{karte}{akzent}{8.26}{\typstolper}{\meta{1}{10 min}{Band 8, Kap.\ 9}}
\kartentitel{Fixpunkte ohne Kompaktheit}
\worum{Brouwer gilt für die abgeschlossene Kugel --- nicht für die offene, nicht für den Kreisring, nicht für die Sphäre. Jede Voraussetzung lässt sich durch ein Gegenbeispiel rechtfertigen.}
\methoden{Gegenbeispiel, Brouwerscher Fixpunktsatz, Kompaktheit}
Geben Sie jeweils eine stetige Selbstabbildung ohne Fixpunkt an:
\begin{enumerate}[label=(\alph*),leftmargin=1.6em,itemsep=2pt]
\item der offenen Einheitskreisscheibe $\{|x| < 1\}$;
\item von $\mathbb R^{2}$;
\item des Kreisrings $\{1\le|x|\le2\}$;
\item der Sphäre $S^{2}$;
\item des Raums $[0, 1]\cup[2, 3]$.
\end{enumerate}
Welche Voraussetzung von Brouwer ist jeweils verletzt --- Kompaktheit, Konvexität (allgemeiner: Zusammenziehbarkeit), Zusammenhang?
\end{karte}

%%%%%%%%%%% Kapitel 10
\stupser{tief}{Nichtorientierbar genau dann, wenn ein Buchstabe zweimal mit demselben Exponenten vorkommt. $\chi = E - K + 1$ mit $K$ = Zahl der verschiedenen Buchstaben; $E$ durch Verfolgen der Ecken.}{22pt}
\begin{karte}{tief}{8.27}{\typtypisch}{\meta{2}{15 min}{Band 8, Kap.\ 10}}
\kartentitel{Kantenwörter lesen}
\worum{Jede geschlossene Fläche ist ein Vieleck mit paarweise verklebten Kanten, beschrieben durch ein Wort. Aus dem Wort liest man Orientierbarkeit und $\chi$ ab --- und damit nach dem Klassifikationssatz die Fläche.}
\methoden{Kantenwort, Euler-Charakteristik, Orientierbarkeit, Klassifikation der Flächen, Verklebung verfolgen}
Bestimmen Sie für jedes Wort die Zahl der Ecken nach dem Verkleben, $\chi$, die Orientierbarkeit und die Fläche:
\begin{enumerate}[label=(\alph*),leftmargin=1.6em,itemsep=2pt]
\item $aba^{-1}b^{-1}$
\item $aa$
\item $aabb$
\item $abca^{-1}b^{-1}c^{-1}$
\item $abcb^{-1}ac$
\item $aba^{-1}b^{-1}cdc^{-1}d^{-1}$
\item Welche Flächen haben $\chi = 0$? Welche $\chi = -1$? Warum gibt es keine orientierbare Fläche mit $\chi = -1$?
\end{enumerate}
\end{karte}

\stupser{gruen}{Eine Kantensorte ist gleichsinnig verklebt (das gibt einen Zylinder), die andere gegensinnig: Die beiden Zylinderränder werden verdreht verklebt --- Kleinsche Flasche.}{22pt}
\begin{karte}{gruen}{8.28}{\typwerk}{\meta{2}{15 min}{Band 8, Kap.\ 10, Aufgabe 10.1}}
\kartentitel{Das Wort $abab^{-1}$}
\worum{Ein Quadrat, ein Wort, eine Fläche: Welche? Die Aufgabe führt die Rechnung vollständig durch --- Ecken verfolgen, $\chi$, Orientierbarkeit, Fundamentalgruppe --- und vergleicht mit dem Torus $aba^{-1}b^{-1}$.}
\methoden{Kantenwort, Euler-Charakteristik, Orientierbarkeit, Fundamentalgruppe, Verklebung verfolgen}
\begin{enumerate}[label=(\alph*),leftmargin=1.6em,itemsep=2pt]
\item Zeichnen Sie das Quadrat mit den Kantenpfeilen und verfolgen Sie, welche Ecken identifiziert werden.
\item Berechnen Sie $\chi$ und entscheiden Sie die Orientierbarkeit. Welche Fläche ist es?
\item Geben Sie $\pi_1$ als Präsentation an (Erzeuger: die Kanten, Relation: das Wort) und zeigen Sie, dass die Gruppe nicht abelsch ist.
\item Zeigen Sie, dass die Fläche sich auch als Verklebung zweier Möbiusbänder längs ihrer Ränder schreiben lässt, und prüfen Sie $\chi$ auf diesem Weg.
\end{enumerate}
\end{karte}

%%%%%%%%%%% Kapitel 11
\stupser{gruen}{$C_2 = \mathbb Z^{4}\to C_1 = \mathbb Z^{6}\to C_0 = \mathbb Z^{4}$. Ränge der Randmatrizen: $\operatorname{rg}\partial_2 = 3$, $\operatorname{rg}\partial_1 = 3$. Dann $b_2 = 4 - 3$, $b_1 = (6 - 3) - 3$, $b_0 = 4 - 3$.}{22pt}
\begin{karte}{gruen}{8.29}{\typwerk}{\meta{2}{20 min}{Band 8, Kap.\ 11, Aufgabe 11.1}}
\kartentitel{Homologie des Tetraederrands}
\worum{Die simpliziale Homologie ist lineare Algebra über $\mathbb Z$: Kettengruppen, Randabbildungen, Kerne modulo Bilder. Die Sphäre als Tetraederoberfläche ist das kleinste Beispiel, an dem alles vorkommt.}
\methoden{Simpliziale Homologie, Randoperator, Rang, Lineare Gleichungssysteme, Euler-Charakteristik}
Ecken $v_0, v_1, v_2, v_3$, alle sechs Kanten $[v_iv_j]$, alle vier Dreiecke $[v_iv_jv_k]$, kein Volumen.
\begin{enumerate}[label=(\alph*),leftmargin=1.6em,itemsep=2pt]
\item Schreiben Sie die Randmatrizen $\partial_1\colon C_1\to C_0$ ($4\times6$) und $\partial_2\colon C_2\to C_1$ ($6\times4$) auf und prüfen Sie $\partial_1\partial_2 = 0$.
\item Bestimmen Sie $H_0 = C_0/\operatorname{im}\partial_1$.
\item Bestimmen Sie $\ker\partial_1$ und $\operatorname{im}\partial_2$ und folgern Sie $H_1 = 0$.
\item Bestimmen Sie $H_2 = \ker\partial_2$ und einen Erzeuger (die \enquote{Fundamentalklasse}).
\item Prüfen Sie $\chi = b_0 - b_1 + b_2 = E - K + F$.
\end{enumerate}
\end{karte}

\stupser{tief}{Mayer--Vietoris für $X = A\cup B$: $\dots\to H_1(A\cap B)\to H_1(A)\oplus H_1(B)\to H_1(X)\to H_0(A\cap B)\to\dots$ Für die Acht: $A\cap B$ zusammenziehbar, also $H_1(X) = H_1(A)\oplus H_1(B)$.}{22pt}
\begin{karte}{tief}{8.30}{\typtypisch}{\meta{2}{20 min}{Band 8, Kap.\ 11}}
\kartentitel{Betti-Zahlen bestimmen}
\worum{Betti-Zahlen zählt man auf drei Wegen: direkt aus einer Triangulierung, mit Mayer--Vietoris aus Teilstücken, oder --- für Flächen --- aus $\chi$ und Orientierbarkeit. Die Aufgabe übt alle drei am Kreis, an der Acht und am Torus.}
\methoden{Simpliziale Homologie, Mayer-Vietoris, Betti-Zahlen, Euler-Charakteristik, Quadratmodell}
\begin{enumerate}[label=(\alph*),leftmargin=1.6em,itemsep=2pt]
\item $S^{1}$ als Rand eines Dreiecks ($3$ Ecken, $3$ Kanten): Bestimmen Sie $H_0$ und $H_1$ direkt.
\item $S^{1}\vee S^{1}$ (Acht): Bestimmen Sie $H_1$ mit Mayer--Vietoris. Welche Rolle spielt, dass $A\cap B$ zusammenziehbar ist?
\item Torus: Das Quadratmodell ist ein Zellkomplex mit einer Ecke, zwei Kanten $a, b$, einer Fläche $F$ mit $\partial F = a + b - a - b$. Berechnen Sie $H_0, H_1, H_2$.
\item Für eine geschlossene orientierbare Fläche gilt $b_0 = b_2 = 1$. Bestimmen Sie $b_1$ für die Fläche vom Geschlecht $g$ aus $\chi = 2 - 2g$. Was ändert sich bei nichtorientierbaren Flächen ($b_2 = 0$, Torsion in $H_1$) --- etwa bei der Kleinschen Flasche?
\end{enumerate}
\end{karte}

%%%%%%%%%%% Kapitel 12
\stupser{gruen}{Invarianz: Bei jedem Reidemeister-Zug lässt sich eine Färbung des alten Diagramms eindeutig auf das neue fortsetzen. Beim Zug II etwa bekommt der neue kleine Bogen die Farbe $2a - b$.}{22pt}
\begin{karte}{gruen}{8.31}{\typwerk}{\meta{2}{15 min}{Band 8, Kap.\ 12, Aufgabe 12.1}}
\kartentitel{Dreifärbbarkeit}
\worum{Ein Diagramm heißt dreifärbbar, wenn man die Bögen so mit drei Farben färben kann, dass an jeder Kreuzung alle drei Farben oder nur eine vorkommen --- und mindestens zwei Farben insgesamt. Es ist die einfachste Invariante, die den Kleeblattknoten vom Unknoten unterscheidet.}
\methoden{Knoteninvariante, Reidemeister-Züge, Dreifärbbarkeit, Fallunterscheidung}
\begin{enumerate}[label=(\alph*),leftmargin=1.6em,itemsep=2pt]
\item Zeigen Sie, dass das Standarddiagramm des Kleeblatts dreifärbbar ist, und dass das kreuzungsfreie Diagramm des Unknotens es nicht ist.
\item Zeigen Sie, dass Dreifärbbarkeit unter den drei Reidemeister-Zügen erhalten bleibt (für Zug III genügt es, die Fälle \enquote{alle drei Bögen gleich} und \enquote{alle drei verschieden} am oberen Bogen zu betrachten).
\item Folgern Sie: Das Kleeblatt ist nicht der Unknoten.
\item Zeigen Sie: Die Farben $\{0, 1, 2\}$ mit der Kreuzungsregel $2\cdot(\text{Überbogen})\equiv\text{Unterbogen}_1 + \text{Unterbogen}_2\pmod3$ sind genau die Dreifärbungen. Wie viele Dreifärbungen hat das Kleeblatt (einfarbige mitgezählt)?
\end{enumerate}
\end{karte}

\stupser{tief}{Die Kreuzungsgleichungen modulo $3$ haben nur die einfarbigen Lösungen. Modulo $5$ gibt es $25$ Lösungen, also $20$ echte Fünffärbungen.}{22pt}
\begin{karte}{tief}{8.32}{\typtypisch}{\meta{3}{20 min}{Band 8, Kap.\ 12}}
\kartentitel{Der Achterknoten}
\worum{Der Achterknoten hat vier Kreuzungen und ist nicht dreifärbbar --- ist er also der Unknoten? Nein: Man muss die Farbenzahl wechseln. Die $p$-Färbungen (Fox) sind ein lineares Gleichungssystem modulo $p$, und welche $p$ funktionieren, verrät die Determinante des Knotens.}
\methoden{Knoteninvariante, Lineare Gleichungssysteme, Modulare Arithmetik, Determinante, Rang}
Das Standarddiagramm des Achterknotens hat vier Bögen $1, 2, 3, 4$ und vier Kreuzungen; an jeder liegt ein Bogen oben und zwei Bögen enden unten:
\begin{center}\small
\begin{tabular}{c|cccc}
Kreuzung & I & II & III & IV\\ \hline
Überbogen & $1$ & $3$ & $4$ & $2$\\
Unterbögen & $2, 4$ & $1, 2$ & $3, 1$ & $4, 3$
\end{tabular}
\end{center}
Eine $p$-Färbung ordnet jedem Bogen eine Zahl in $\mathbb Z/p$ zu, so dass an jeder Kreuzung $2\cdot\text{über}\equiv\text{unter}_1 + \text{unter}_2\pmod p$.
\begin{enumerate}[label=(\alph*),leftmargin=1.6em,itemsep=2pt]
\item Schreiben Sie die vier Gleichungen als Matrix $M$ (Zeilen: Kreuzungen, Spalten: Bögen). Zeigen Sie, dass jede Zeilen- und jede Spaltensumme null ist.
\item Zeigen Sie, dass es modulo $3$ nur einfarbige Lösungen gibt. Der Achterknoten ist also nicht dreifärbbar.
\item Finden Sie eine echte Fünffärbung. Wie viele gibt es?
\item Streichen Sie eine Zeile und eine Spalte von $M$ und berechnen Sie die Determinante des Rests. Für welche Primzahlen $p$ ist der Knoten $p$-färbbar? Vergleichen Sie mit dem Kleeblatt (Determinante $3$).
\item Folgern Sie, dass der Achterknoten weder der Unknoten noch das Kleeblatt ist.
\end{enumerate}
\end{karte}

%%%%%%%%%%% Kapitel 13
\stupser{gruen}{Liegt der Torus flach auf dem Tisch (Achse senkrecht), ist die Höhe $f = \sin u$ auf dem obersten und untersten Kreis konstant --- kritische Kreise statt Punkte. Steht er auf der Seite, ist $f = (2 + \cos u)\cos v$.}{22pt}
\begin{karte}{gruen}{8.33}{\typwerk}{\meta{2}{20 min}{Band 8, Kap.\ 13, Aufgabe 13.1}}
\kartentitel{Höhenfunktionen auf dem Torus}
\worum{Morse-Theorie liest die Topologie einer Mannigfaltigkeit an den kritischen Punkten einer Funktion ab: $c_0 - c_1 + c_2 = \chi$. Die Aufgabe vergleicht den stehenden Torus (vier isolierte kritische Punkte) mit dem liegenden (zwei kritische Kreise).}
\methoden{Kritische Punkte, Hesse-Matrix, Morse-Index, Euler-Charakteristik, Partielle Ableitungen}
Torus mit senkrechter Symmetrieachse: $X(u, v) = \bigl((2 + \cos u)\cos v, (2 + \cos u)\sin v, \sin u\bigr)$.
\begin{enumerate}[label=(\alph*),leftmargin=1.6em,itemsep=2pt]
\item Liegender Torus (flach auf dem Tisch), $f = z = \sin u$: Bestimmen Sie die kritischen Punkte. Warum ist $f$ keine Morse-Funktion?
\item Stehender Torus (auf der Seite), $f = x = (2 + \cos u)\cos v$: Berechnen Sie $f_u$, $f_v$ und die vier kritischen Punkte.
\item Bestimmen Sie dort die Hesse-Matrix und die Indizes (Minimum $0$, Sattel $1$, Maximum $2$) und prüfen Sie $c_0 - c_1 + c_2 = \chi(T^{2})$.
\item Was bedeutet \enquote{Morse-Funktionen liegen dicht}, und wie \enquote{repariert} ein kleines Kippen den liegenden Torus?
\end{enumerate}
\end{karte}

%%%%%%%%%%% Kapitel 14
\stupser{gruen}{Abstände im regelmäßigen Achteck: $2\sin22{,}5^{\circ} = 0{,}765$, $\sqrt2 = 1{,}414$, $2\sin67{,}5^{\circ} = 1{,}848$, $2$. Bei jedem Wert kommen Kanten hinzu --- und manchmal Dreiecke.}{22pt}
\begin{karte}{gruen}{8.34}{\typwerk}{\meta{2}{20 min}{Band 8, Kap.\ 14, Aufgabe 14.1}}
\kartentitel{Ein Barcode von Hand}
\worum{Persistente Homologie: Man verbindet Punkte, deren Abstand höchstens $\varepsilon$ ist (Vietoris--Rips-Komplex), lässt $\varepsilon$ wachsen und protokolliert, wann Löcher entstehen und verschwinden. Acht Punkte auf einem Kreis sind klein genug für die Hand.}
\methoden{Vietoris-Rips-Komplex, Betti-Zahlen, Euler-Charakteristik, Filtrierung, Trigonometrie}
Die Punkte $p_k = (\cos\frac{2\pi k}8, \sin\frac{2\pi k}8)$, $k = 0, \dots, 7$.
\begin{enumerate}[label=(\alph*),leftmargin=1.6em,itemsep=2pt]
\item Berechnen Sie die vier vorkommenden Abstände.
\item Beschreiben Sie den Rips-Komplex für $\varepsilon$ knapp unter und knapp über jedem der vier Werte: Zahl der Ecken, Kanten, Dreiecke (höhere Simplizes nur qualitativ), $b_0$ und $b_1$. Prüfen Sie mit $\chi$, wo es geht.
\item Bei welchem $\varepsilon$ wird das zentrale Loch geboren, bei welchem stirbt es? Zeichnen Sie den Barcode für $H_0$ und $H_1$.
\item Was passiert bei $\varepsilon = 2\sin67{,}5^{\circ}$ in höheren Dimensionen --- welcher bekannte Raum ist der Rips-Komplex dort? (Jeder Punkt ist mit allen außer seinem Gegenüber verbunden.)
\end{enumerate}
\end{karte}
